\documentclass[11pt,a4paper]{article}
\usepackage[utf8]{inputenc}
\usepackage[T1]{fontenc}
\usepackage{geometry}
\geometry{margin=1in}
\usepackage{hyperref}
\usepackage{enumitem}
\usepackage{booktabs}
\usepackage{longtable}
\usepackage{xcolor}
\usepackage{titlesec}
\usepackage{amsmath,amssymb}

\titleformat{\section}{\Large\bfseries\color{blue!70!black}}{}{0em}{}
\titleformat{\subsection}{\large\bfseries\color{blue!50!black}}{}{0em}{}
\titleformat{\subsubsection}{\normalsize\bfseries\color{blue!30!black}}{}{0em}{}

\title{\textbf{Replication Plan}\\[0.5em]
\large Physics-Based Hybrid Machine Learning for Critical Heat Flux Prediction with Uncertainty Quantification\\[0.3em]
\normalsize OSTI ID: 2571909}
\author{Auto-generated Replication Blueprint}
\date{\today}

\begin{document}
\maketitle
\tableofcontents
\newpage

%---------------------------------------------------------------------
\section{Paper Overview}
%---------------------------------------------------------------------
This paper develops a physics-based hybrid machine learning approach for predicting critical heat flux (CHF) in nuclear reactor thermal-hydraulics. The hybrid model combines established empirical CHF correlations (e.g., Groeneveld look-up table, W-3 correlation) as physics-based features with data-driven models (deep neural networks, Bayesian neural networks, deep Gaussian processes) to improve prediction accuracy. Uncertainty quantification (UQ) is provided through Bayesian approaches. The model is trained and validated on the NRC/EPRI CHF database.

%---------------------------------------------------------------------
\section{Detailed Workflow}
%---------------------------------------------------------------------

\subsection{Phase 1: Data Acquisition and Preprocessing}
\begin{enumerate}[label=\textbf{1.\arabic*}]
  \item \textbf{Obtain the CHF database.}
    \begin{itemize}
      \item NRC/EPRI CHF database or the publicly available subset.
      \item Features: pressure, mass flux, quality, heated length, hydraulic diameter, etc.
      \item Target: critical heat flux (CHF) value.
    \end{itemize}
  \item \textbf{Filter and clean data.}
    \begin{itemize}
      \item Remove outliers and entries with missing values.
      \item Apply parameter range restrictions as in the paper.
    \end{itemize}
  \item \textbf{Feature engineering.}
    \begin{itemize}
      \item Compute physics-based features: correlation predictions (W-3, look-up table, etc.).
      \item Normalize/standardize inputs.
    \end{itemize}
  \item \textbf{Split data} into train/validation/test sets (ratios as specified).
\end{enumerate}

\subsection{Phase 2: Baseline Physics Correlations}
\begin{enumerate}[label=\textbf{2.\arabic*}]
  \item \textbf{Implement CHF correlations.}
    \begin{itemize}
      \item Groeneveld 2006 look-up table.
      \item W-3 correlation.
      \item Other correlations referenced in the paper.
    \end{itemize}
  \item \textbf{Evaluate correlation performance} on the test set (RMSE, $R^2$, MAPE).
\end{enumerate}

\subsection{Phase 3: Hybrid ML Model Building}
\begin{enumerate}[label=\textbf{3.\arabic*}]
  \item \textbf{Deep Neural Network (DNN).}
    \begin{itemize}
      \item Input: physical features + correlation predictions.
      \item Architecture: fully connected, 3--5 hidden layers, ReLU activation.
      \item Train with Adam optimizer, MSE loss.
    \end{itemize}
  \item \textbf{Bayesian Neural Network (BNN).}
    \begin{itemize}
      \item Variational inference or MC-Dropout for uncertainty estimation.
      \item Output: predictive mean and variance.
    \end{itemize}
  \item \textbf{Deep Gaussian Process (DGP).}
    \begin{itemize}
      \item Stacked GP layers for capturing non-stationary relationships.
      \item Inference via variational methods (e.g., doubly stochastic VI).
    \end{itemize}
\end{enumerate}

\subsection{Phase 4: Uncertainty Quantification}
\begin{enumerate}[label=\textbf{4.\arabic*}]
  \item \textbf{Compute predictive intervals} (e.g., 95\% credible intervals).
  \item \textbf{Calibration analysis}: compare predicted vs.\ observed coverage.
  \item \textbf{Decompose uncertainty}: aleatoric vs.\ epistemic components.
\end{enumerate}

\subsection{Phase 5: Validation}
\begin{enumerate}[label=\textbf{5.\arabic*}]
  \item \textbf{Compare hybrid vs.\ pure-ML vs.\ pure-correlation} approaches.
  \item \textbf{Reproduce error metrics} (RMSE, MAE, $R^2$) from paper tables.
  \item \textbf{Reproduce parity plots, residual distributions, UQ coverage plots.}
\end{enumerate}

%---------------------------------------------------------------------
\section{Datasets}
%---------------------------------------------------------------------

\begin{longtable}{p{4.5cm} p{3.5cm} p{6cm}}
\toprule
\textbf{Dataset / Source} & \textbf{Type} & \textbf{Location / Notes} \\
\midrule
\endfirsthead
\toprule
\textbf{Dataset / Source} & \textbf{Type} & \textbf{Location / Notes} \\
\midrule
\endhead
NRC/EPRI CHF database & Experimental & Contact NRC or EPRI for access; some subsets available in literature \\
\midrule
Groeneveld 2006 look-up table & Correlation table & Published in: Groeneveld et al., Nuclear Engineering and Design (2007) \\
\midrule
Paper's data splits & Specified & Training/test splits described in paper methodology \\
\bottomrule
\end{longtable}

%---------------------------------------------------------------------
\section{Models, Tools, and Software}
%---------------------------------------------------------------------

\begin{longtable}{p{4.5cm} p{2.5cm} p{6.5cm}}
\toprule
\textbf{Tool / Library} & \textbf{Version} & \textbf{Purpose / URL} \\
\midrule
\endfirsthead
\toprule
\textbf{Tool / Library} & \textbf{Version} & \textbf{Purpose / URL} \\
\midrule
\endhead
PyTorch & $\geq 1.10$ & DNN and BNN implementation \\
\midrule
GPyTorch & $\geq 1.9$ & Deep Gaussian process \newline \url{https://gpytorch.ai/} \\
\midrule
Pyro & $\geq 1.8$ & Probabilistic programming for BNN \newline \url{https://pyro.ai/} \\
\midrule
scikit-learn & $\geq 1.0$ & Data preprocessing, metrics, baseline models \\
\midrule
NumPy / SciPy & $\geq 1.21$ & Numerical operations \\
\midrule
Pandas & $\geq 1.3$ & Data manipulation \\
\midrule
Matplotlib / Seaborn & $\geq 3.4$ & Parity plots, residual plots, UQ visualization \\
\bottomrule
\end{longtable}

\subsection{Hardware Requirements}
\begin{itemize}
  \item GPU recommended for BNN and DGP training ($\geq$ 8\,GB VRAM).
  \item $\geq$ 16\,GB RAM.
\end{itemize}

\end{document}
